% CS 418: Intro to Data Science — Worksheet 16 (16-week08-day2)
% Week 8, Wednesday 2026-10-14
% NOTICE: This worksheet was drafted by an LLM coding system from the
% instructor's choice of checkpoints. It is maintained, reviewed, and vetted by humans.
% Compile with: latexmk -pdf worksheet.tex

\documentclass[11pt]{article}

% ── Packages ──────────────────────────────────────────────────────────────────
\usepackage[margin=0.65in, headheight=14pt]{geometry}
\usepackage{amsmath, amssymb}
\usepackage{ragged2e}
% Never break a word across lines, including at hyphens ("Sans-Serif").
\hyphenpenalty=10000
\exhyphenpenalty=10000
\usepackage{enumitem}
\usepackage{booktabs}
\usepackage{graphicx}
\usepackage{hyperref}
\usepackage{xcolor}
\usepackage{fancyhdr}
\usepackage{mdframed}
\usepackage{array}

% ── Page style ────────────────────────────────────────────────────────────────
\pagestyle{fancy}
\fancyhf{}
\lhead{\textbf{CS 418: Intro to Data Science}}
\rhead{UIC, Fall 2026}
\cfoot{\thepage}
\renewcommand{\headrulewidth}{1pt}
\setlength{\headheight}{12pt}
\setlength{\headsep}{0.4cm}

% ── Hyperlinks ────────────────────────────────────────────────────────────────
\hypersetup{colorlinks=false, hidelinks}

% ── Convenience environments ──────────────────────────────────────────────────
% Usage: \blankspace{6cm} — blank area for hand-written work, no rules
\newcommand{\blankspace}[1]{%
  \vspace{0.3em}
  \begin{minipage}[t][#1][t]{\linewidth}
  \end{minipage}
  \vspace{0.3em}
}

% Usage: \guidedopen{title}{guided content}{guided workspace height}{open-ended workspace height}
% Left: a titled, structured prompt with its own blank workspace. Right: a fully blank workspace for open-ended exploration of the same topic (no prompt text), sized so the two columns land at roughly the same total height. Neither workspace is ruled.
\newcommand{\guidedopen}[4]{%
  \noindent\textbf{#1}\\[0.18em]
  \noindent
  \begin{minipage}[t]{0.485\textwidth}
    #2
    \blankspace{#3}
  \end{minipage}
  \hfill
  \begin{minipage}[t]{0.485\textwidth}
    \blankspace{#4}
  \end{minipage}
  \\[0.4em]
}

% Usage: \panelbox{width}{height}{caption} — a framed drawing panel with `caption` centered beneath it. An empty caption leaves the panel bare, with nothing below the frame at all.
\newcommand{\panelbox}[3]{%
  \begin{minipage}[t]{\dimexpr#1+2\fboxsep+2\fboxrule\relax}
    \framebox[#1][c]{\rule{0pt}{#2}}%
    \if\relax\detokenize{#3}\relax\else
      \\[0.25em]{\centering\small #3\par}%
    \fi
  \end{minipage}%
}

% Usage: \cvrow{label} — one row of the cross-validation grid: a label, then five empty fold cells.
\newcommand{\cvrow}[1]{\small #1 & \rule{0pt}{0.6cm} & & & & \\ \cline{2-6}}

% ══════════════════════════════════════════════════════════════════════════════
\begin{document}
\RaggedRight

% ── Title block: topic left, info box right ───────────────────────────────────
\noindent
\begin{minipage}[t]{0.55\textwidth}
  \vspace{0pt}
  {\LARGE \textbf{Worksheet \#\,16}} \\[0.18em]
  {\large \textit{Week 8, Day 2: Intro to classification}}
\end{minipage}
\hfill
\begin{minipage}[t]{0.40\textwidth}
  \vspace{0pt}
  \small
  Name: \underline{\hspace{4.5cm}} \\[0.35em]
  NetID: \underline{\hspace{4.3cm}} \\[0.35em]
  Date: \underline{\hspace{4.4cm}} \\
\end{minipage}

\vspace{0.6em}
\noindent\rule{\linewidth}{0.6pt}
\vspace{0.1em}

\noindent\underline{\textbf{Guided checkpoints}}\\[0.2em]
\noindent{\small For full credit, fill in the checkpoints. If you skip a checkpoint, you must include your own notes for full credit.}
\vspace{0.8em}

% ── 1. Four steps to modeling ────────────────────────────────────────────────
\guidedopen{1. Four Steps to Modeling}{
Write down the four main steps for modeling, ideally from memory!

\vspace{0.5em}
1.
\vspace{0.7cm}

2.
\vspace{0.7cm}

3.
\vspace{0.7cm}

4.
}{0.7cm}{4.4cm}

\vfill

% ── 2. LINE ──────────────────────────────────────────────────────────────────
\guidedopen{2. LINE Criteria}{
Write down the LINE criteria for applying linear regression to a dataset, ideally from memory.

\vspace{0.5em}
L:
\vspace{0.7cm}

I:
\vspace{0.7cm}

N:
\vspace{0.7cm}

E:
}{0.7cm}{4.8cm}

\vfill

% ── 3. Regression vs. classification ─────────────────────────────────────────
\guidedopen{3. Regression vs. Classification}{
Sketch a regression model on the left and a classifier on the right. Write what each model predicts.

\vspace{0.45em}

\noindent
\panelbox{3.7cm}{3.2cm}{regression}\hfill
\panelbox{3.7cm}{3.2cm}{classification}

\vspace{0.5em}
\noindent\makebox[3.9cm][l]{Predicts:}\hfill\makebox[3.9cm][l]{Predicts:}
}{0.6cm}{6.0cm}

\vfill

\newpage

% ── 4. Three sets ────────────────────────────────────────────────────────────
\guidedopen{4. Three Sets}{
Split the bar into training, validation, and test sets. Label each part with its share of the data and write what each set is used for.

\vspace{0.45em}

\noindent
\panelbox{7.9cm}{0.9cm}{}

\vspace{0.6em}
Training:
\vspace{0.7cm}

Validation:
\vspace{0.7cm}

Test:
}{0.7cm}{5.0cm}

\vfill

% ── 5. Cross-validation ──────────────────────────────────────────────────────
\guidedopen{5. Cross-Validation}{
Fill in each cell with V (validation) or T (training) for 5-fold cross-validation.

\vspace{0.5em}

\noindent
\begin{tabular}{r|*{5}{>{\centering\arraybackslash}p{1.0cm}|}}
  \multicolumn{1}{r}{} & \multicolumn{1}{c}{\small Fold 1} & \multicolumn{1}{c}{\small Fold 2} & \multicolumn{1}{c}{\small Fold 3} & \multicolumn{1}{c}{\small Fold 4} & \multicolumn{1}{c}{\small Fold 5} \\ \cline{2-6}
  \cvrow{Split 1}
  \cvrow{Split 2}
  \cvrow{Split 3}
  \cvrow{Split 4}
  \cvrow{Split 5}
\end{tabular}
}{0.4cm}{4.6cm}

\vfill

% ── 6. Project ideas ─────────────────────────────────────────────────────────
\guidedopen{6. Models for Your Project}{
Write down two ideas for using a regression model and/or a classifier in your project.

\vspace{0.5em}
Idea 1:
\vspace{3.0cm}

Idea 2:
}{3.0cm}{7.4cm}

\vfill

\end{document}
